\section{模型的检验与灵敏度分析}

\subsection{数值可行性检验}

对四个问题的全部执行方案逐区间回代检验，结果见表 12。能量平衡按
$x_t+PV_t+d_t+E_t=L_t+c_t+w_t$ 核算，储电量与充放电量按式~\eqref{eq:p1soc} 递推。

\threelinetable{表 12\quad 数值可行性检验的最大违反量}{lrrrrr}
{检验项 & 问题二 & 问题三 & 问题 4-2 & 问题 4-3 & 阈值}
{能量平衡残差 / kWh & $4.6\times10^{-13}$ & $2.3\times10^{-13}$ & $2.3\times10^{-13}$ & $2.3\times10^{-13}$ & $10^{-6}$ \\
 储电量越界 / kWh & $0$ & $3.0\times10^{-12}$ & $0$ & $3.6\times10^{-12}$ & $10^{-6}$ \\
 充放电功率越界 / kWh & $0$ & $0$ & $0$ & $1.8\times10^{-12}$ & $10^{-6}$ \\
 售电（负购电量）/ kWh & $0$ & $0$ & $0$ & $0$ & $0$}

另外核对过：五个结果文件的分段值之和与“全天购电量”列完全一致（最大偏差 $10^{-11}$，来自四位小数舍入）；\texttt{result2.xlsx} 的充放电量表共 $2005=1+334\times6$ 行，覆盖 2 月至 12 月全部 334 天，无遗漏。

\subsection{因果性检验}

对问题二至问题四的每一天，程序断言其使用的预测只由第 $d-1$ 天及更早的数据构成；问题三的调整只使用发布时刻不晚于该时刻的预报。以式~\eqref{eq:fc} 与式~\eqref{eq:p3expand} 的实现为例，第 $d$ 天的输入切片上界恒为 $d$（开区间），不存在读取当天未来的可能。

\subsection{参数灵敏度}

以问题二全年总费用 $17026400$ 元为基准，逐项扰动关键参数，结果见表 13。

\threelinetable{表 13\quad 关键参数灵敏度（问题二全年，单位：元）}{llrr}
{参数 & 取值 & 总费用 & 相对基准}
{日前预测窗口 & 前 1 日 & 20866511 & $+3840111$ \\
              & 前 3 日 & 20544768 & $+3518368$ \\
              & \textbf{前 7 日} & \textbf{17026400} & $\mathbf{0}$ \\
              & 前 14 日 & 17150375 & $+123975$ \\
              & 前 30 日 & 17326664 & $+300264$ \\
 安全裕度分位点 & 0.00（不加裕度） & 23959895 & $+6933495$ \\
              & 0.70 & 17163260 & $+136860$ \\
              & \textbf{0.80} & \textbf{17026400} & $\mathbf{0}$ \\
              & 0.90 & 17258273 & $+231873$ \\
 误差样本窗口 & 前 7 日 & 17284369 & $+257969$ \\
              & 前 28 日 & 17026400 & $0$ \\
              & 前 120 日 & 16986934 & $-39467$ \\
 收尾电量约束 & 日循环（本文） & 17026400 & $0$ \\
              & 收尾自由 & 17053078 & $+26678$ \\
 充电效率 $\eta_c$ & 0.85 & 17193129 & $+166728$ \\
              & \textbf{0.90} & \textbf{17026400} & $\mathbf{0}$ \\
              & 1.00 & 16700941 & $-325459$}

三点结论：

\textbf{（1）7 日预测窗口是一个真实的极值点。}窗口从 1 日加长到 7 日，费用从 $2086.7$ 万元迅速降到 $1702.6$ 万元（降幅 $18.4\%$）；再继续加长到 30 日反而回升到 $1732.7$ 万元。原因是短窗口噪声大、误差分位数估不准，长窗口则把季节性的负载水平抹平了。

\textbf{（2）安全裕度分位点对费用极其敏感，且最优值可被理论预测。}不加裕度要多花 $693$ 万元（$40.7\%$），分位点每偏离 $0.1$ 约损失 $13\sim23$ 万元。表 14 进一步验证了式~\eqref{eq:newsvendor}：当紧急购电价改为 $m$ 倍时，理论最优分位点 $q=1-1/m$ 始终优于固定取 $0.8$。

\threelinetable{表 14\quad 紧急购电价比率 $m$ 与最优安全裕度}{crrr}
{$m$ & 理论 $q=1-1/m$ & 按理论 $q$ 的费用 / 元 & 固定 $q=0.8$ 的费用 / 元}
{2.0  & 0.500 & 15604004 & 16232095 \\
 3.0  & 0.667 & 16299627 & 16496864 \\
 5.0  & 0.800 & 17026400 & 17026400 \\
 8.0  & 0.875 & 17648981 & 17820705 \\
 10.0 & 0.900 & 17946676 & 18350242}

\textbf{（3）效率口径与收尾约束的影响是可控的。}充电效率在 $0.85\sim1.00$ 之间变动时，全年费用变化在 $-32.5$ 万到 $+16.7$ 万元之间（不超过 $2\%$），说明“充放电效率 $90\%$”这一歧义的量级有限。收尾约束方面，允许储电量跨日自由漂移反而贵 $2.7$ 万元——因为收尾电量自由时线性规划会把储能一路充到 $10800\,\mathrm{kWh}$ 的上限并长期停在上面（把残值折现的退化行为），所以日循环口径既更贴近题面，也确实更省。
